% In compendium/laws-of-harmony-and-dynamics/law-of-wobble-periodicity.tex
\subsection*{The Law of Wobble Periodicity}
\hrule
\vspace{0.3cm}
\GetLaw{LawOfWobblePeriodicity}
\vspace{0.5cm}

\subsubsection*{Conceptual Overview}
This law asserts that the fundamental 'dissonance' or imbalance in the framework is not a chaotic element but a perfectly ordered, periodic rotation. The characteristic recurrence proves that this 'wobble' is not random noise but is instead a direct and pure expression of the Golden Algebra itself.

\subsubsection*{Proof}
The proof verifies that the dynamics described by the law's characteristic recurrence coefficients constitute a pure, non-decaying rotation. This is done by constructing the characteristic polynomial from these coefficients and proving its roots (the system's eigenvalues) lie on the unit circle.

\subsubsection*{Formal Proof}
\begin{enumerate}
    \item The law states the dynamics are governed by a recurrence with coefficients native to the Golden Algebra:
    $$ c_1 = \Phi \quad \text{and} \quad c_2 = 2(T+K) $$
    Using the proven identity $T+K = -1/2$ (Property 14), we find $c_2 = -1$.
    
    \item The characteristic polynomial for such a recurrence is $x^2 - c_1x - c_2 = 0$. Substituting the coefficients gives:
    $$ x^2 - \Phi x + 1 = 0 $$
    
    \item The roots of this polynomial, $\lambda$, are the eigenvalues that govern the system's dynamics. For a quadratic equation $ax^2+bx+c=0$, the product of the roots is $c/a$. Here, the product of the eigenvalues is $\lambda_1 \lambda_2 = 1/1 = 1$.
    
    \item If the roots are complex conjugates, such that $\lambda_2 = \overline{\lambda_1}$, then their product is the squared magnitude: $\lambda_1 \overline{\lambda_1} = |\lambda_1|^2 = 1$. This implies the magnitude itself is $|\lambda_1|=1$.
    
    \item To confirm the roots are complex, we check the discriminant, $\Delta = b^2 - 4ac$:
    $$ \Delta = (-\Phi)^2 - 4(1)(1) = \Phi^2 - 4 $$
    Since $\Phi = (\sqrt{5}-1)/2 \approx 0.618$, $\Phi^2 \approx 0.382$. Thus, $\Delta = \Phi^2 - 4$ is negative.
    
    \item Because the discriminant is negative, the roots are complex conjugates. Because their product is 1, they must lie on the unit circle. This proves the dynamic is a pure, non-decaying rotation. The law is therefore \textbf{proven}.
\end{enumerate}

\subsubsection*{Computational Verification}
The following Python script computationally proves the law. It derives the eigenvalues from the law's coefficients and asserts that their magnitudes are exactly 1.

\begin{lstlisting}[language=Python, caption={Symbolic proof of the Law of Wobble Periodicity.}, label={lst:wobble_periodicity_proof}]
import sympy


def prove_wobble_periodicity():
    """
    Symbolically proves the Law of Wobble Periodicity.

    This script constructs the characteristic polynomial from the coefficients
    given in the law (c1=Phi, c2=2(T+K)), finds its roots (the eigenvalues
    of the dissonance operator), and proves these eigenvalues lie on the
    unit circle, which corresponds to a pure, non-decaying rotation.
    """
    # --- Define symbolic constants ---
    sqrt5 = sympy.sqrt(5)
    T = (sqrt5 - 1) / 4
    K = -(sqrt5 + 1) / 4
    # Golden Conjugate, Phi = -1/2 + sqrt(5)/2
    Phi = (sqrt5 - 1) / 2
    x = sympy.Symbol('x')

    print("=" * 60)
    print("Symbolic Proof of the Law of Wobble Periodicity")
    print("=" * 60)

    # --- Step 1: Define the recurrence coefficients from the Law ---
    c1 = Phi
    c2 = 2 * (T + K)

    print("1. Defining the characteristic recurrence coefficients from the law:")
    print(f"   c1 = Phi = {c1.simplify()}")
    print(f"   c2 = 2*(T+K) = {c2.simplify()}\n")

    # --- Step 2: Construct the characteristic polynomial ---
    # The polynomial is x^2 - c1*x - c2 = 0
    characteristic_poly = x ** 2 - c1 * x - c2
    print("2. Constructing the characteristic polynomial x^2 - c1*x - c2 = 0:")
    print(f"   {characteristic_poly} = 0\n")

    # --- Step 3: Find the eigenvalues (roots) of the polynomial ---
    eigenvalues = sympy.solve(characteristic_poly, x)
    print("3. Solving for the eigenvalues (roots) of the polynomial...")
    print(f"   Eigenvalue 1: {eigenvalues[0]}")
    print(f"   Eigenvalue 2: {eigenvalues[1]}\n")

    # --- Step 4: Verify the eigenvalues lie on the unit circle ---
    # We do this by proving their absolute value (magnitude) is 1.
    print("4. Verifying the eigenvalues lie on the unit circle (magnitude = 1)...")

    is_proven = True
    for i, val in enumerate(eigenvalues):
        magnitude = sympy.simplify(sympy.Abs(val))
        print(f"   Magnitude of Eigenvalue {i + 1}: {magnitude}")
        if magnitude != 1:
            is_proven = False

    assert is_proven, "Proof Failed: An eigenvalue is not on the unit circle."
    print("\n   The magnitude of both eigenvalues is exactly 1.\n")

    print("Conclusion: The eigenvalues of the dissonance operator lie on the")
    print("unit circle, proving the dynamic is a pure, non-decaying rotation.")
    print("The Law of Wobble Periodicity is computationally verified.")
    print("=" * 60)


if __name__ == '__main__':
    prove_wobble_periodicity()
\end{lstlisting}

\paragraph{Verification Output}
\begin{verbatim}
============================================================
Symbolic Proof of the Law of Wobble Periodicity
============================================================
1. Defining the characteristic recurrence coefficients from the law:
    c1 = Phi = -1/2 + sqrt(5)/2
    c2 = 2*(T+K) = -1

2. Constructing the characteristic polynomial x^2 - c1*x - c2 = 0:
    x**2 - x*(-1/2 + sqrt(5)/2) + 1 = 0

3. Solving for the eigenvalues (roots) of the polynomial...
    Eigenvalue 1: -1/4 + sqrt(5)/4 - I*sqrt(2*sqrt(5) + 10)/4
    Eigenvalue 2: -1/4 + sqrt(5)/4 + sqrt(-10 - 2*sqrt(5))/4

4. Verifying the eigenvalues lie on the unit circle (magnitude = 1)...
    Magnitude of Eigenvalue 1: 1
    Magnitude of Eigenvalue 2: 1

    The magnitude of both eigenvalues is exactly 1.

Conclusion: The eigenvalues of the dissonance operator lie on the
unit circle, proving the dynamic is a pure, non-decaying rotation.
The Law of Wobble Periodicity is computationally verified.
============================================================
\end{verbatim}